均匀无粘可压缩流体的密度为 \(\rho>0\)，声速 \(c>0\)。半径 \(a>0\) 的球形声源作均匀小幅径向呼吸运动，边界法向速度是 \(u_r(a,t)=\operatorname{Re}(V_0e^{-i\omega t})\)，\(\omega>0\)。流场在 \(r\ge a\) 无其他声源，只有出射而无入射球面波；速度势满足 \(u_r=\partial_r\phi\)、\(p=-\rho\phi_t\)，且 \(k=\omega/c\)。
(1) 从连续性、线性Euler和声学状态方程建立速度势波动方程，解球对称Helmholtz问题并严格用球面速度边界确定复振幅，求任意 \(r\ge a\) 的声压与速度。
(2) 计算球面比声辐射阻抗 \(Z_{\rm rad}=\widehat p(a)/V_0\)，证明辐射阻抗实部非负，推导其 \(ka\ll1\) 与 \(ka\gg1\) 极限并解释无功项。
(3) 从声强的复幅值表达式计算任意观测球面的总平均辐射功率，核对源面输入功率，推导低频固定体积速度 \(Q_0=4\pi a^2V_0\) 时的辐射幂律。
(4) 比较近场 \(kr\ll1\) 和远场 \(kr\gg1\) 的压强、速度相位及距离标度；计算复声强的虚部、解释其不贡献远场声功，并讨论固定 \(Q_0\) 的 \(a\to0\) 点源极限为什么不能无限保持线性边界速度假设。